% GENERATED FILE. Edit canonical lessons, then run npm run generate:books.
% canonical-source-hash: fnv1a64:69313bd01de659c0
% canonical-lessons: MW-C164-khulo, MW-C164-band, MW-C164-tai, MW-C164-dairy-notice, MW-R164-dairy

\chapter{A Notice at the Dairy}
\label{ch:mw-a2-a-notice-at-the-dairy}
\hlchaptermodality{164}

\begin{chapteropening}
\textbf{By the end of this chapter:} \emph{I can read a dairy's notice of opening days and hours, with a price, and decide when to go.}

Complete MW-R164-dairy: read a 51-word dairy notice and decide when to go for milk, and on which day not to.
\end{chapteropening}

\section[\texorpdfstring{\mw{खुलो} (khulo)}{खुलो (khulo)}]{\mw{खुलो} (khulo) — open}

\label{lesson:MW-C164-khulo}



\begin{quote}
Before the new one: say \emph{to open}, \textbf{\mw{खोलणो}}. Then \emph{a house} and \emph{a shop}.
\end{quote}

\subsection*{You'll want to know: \mw{खुलो}}
\textbf{\mw{खुलो}} — \emph{khulo} — \textquotedblleft{}open\textquotedblright{}.

\begin{quote}
\textbf{\mw{घर} \mw{खुलो} \mw{है।}}
\end{quote}

\begin{quote}
\emph{ghar khulo hai.} — The house is open.
\end{quote}

\begin{cousinweb}
\textbf{\mw{खुलो}} (\emph{khulo}) is inherited: CDIAL lists Old Marwari \emph{khulo}, \textquotedblleft{}open\textquotedblright{}, in the same entry as \textbf{\mw{खोलणो}}, under a reconstructed \emph{khōll}, \textquotedblleft{}to open\textquotedblright{} (CDIAL 3945).
\end{cousinweb}

\begin{grammarlens}[title={What you've built}]
\textbf{\mw{खुलो}} ends like \textbf{\mw{छोरो}} and changes the same way, to agree with what is open: \textbf{\mw{घर} \mw{खुलो} \mw{है।}}, but \textbf{\mw{दुकान} \mw{खुली} \mw{है।}} (\emph{dukān khulī hai.}), the shop is open.
\end{grammarlens}

\subsection*{Your turn}
Say these aloud:

\begin{itemize}
\raggedright
  \item \emph{khulo}
  \item \textbf{\mw{दुकान} \mw{आज} \mw{खुली} \mw{है}?} — \emph{dukān āj khulī hai?} — Is the shop open today?
  \item \textbf{\mw{हां} \mw{सा}, \mw{दुकान} \mw{खुली} \mw{है।}} — \emph{hā̃ sā, dukān khulī hai.} — Yes, the shop is open.
\end{itemize}

\begin{tcolorbox}[breakable,colback=teal!4,colframe=teal!35!black,title={Before you move on}]
What does \textbf{\mw{खुलो}} mean? (\textquotedblleft{}Open\textquotedblright{}.) \emph{The shop is open}? (\textbf{\mw{दुकान} \mw{खुली} \mw{है।}})
\end{tcolorbox}
\section[\texorpdfstring{\mw{बंद} (band)}{बंद (band)}]{\mw{बंद} (band) — closed}

\label{lesson:MW-C164-band}



\begin{quote}
Before the new one: say \emph{the shop is open}. Then \emph{Monday} and \emph{Sunday}: \textbf{\mw{सोमवार}}, \textbf{\mw{रविवार}}.
\end{quote}

\subsection*{You'll want to know: \mw{बंद}}
\textbf{\mw{बंद}} — \emph{band} — \textquotedblleft{}closed\textquotedblright{}.

\begin{quote}
\textbf{\mw{सोमवार} \mw{दुकान} \mw{बंद} \mw{है।}}
\end{quote}

\begin{quote}
\emph{somvār dukān band hai.} — On Monday the shop is closed.
\end{quote}

\begin{cousinweb}
\textbf{\mw{बंद}} (\emph{band}) is a loan, usually traced to Persian \emph{band}, \textquotedblleft{}shut, bound\textquotedblright{}.
\end{cousinweb}

\begin{grammarlens}[title={What you've built}]
\textbf{\mw{बंद}} does not change: \textbf{\mw{घर} \mw{बंद} \mw{है।}}, \textbf{\mw{दुकान} \mw{बंद} \mw{है।}} Its opposite does: \textbf{\mw{खुलो}}, \textbf{\mw{खुली}}.
\end{grammarlens}

\subsection*{Your turn}
Say these aloud:

\begin{itemize}
\raggedright
  \item \emph{band}
  \item \textbf{\mw{दुकान} \mw{बंद} \mw{है} \mw{पण} \mw{घर} \mw{खुलो} \mw{है।}} — \emph{dukān band hai paṇ ghar khulo hai.} — The shop is closed, but the house is open.
  \item \textbf{\mw{दुकान} \mw{आज} \mw{बंद} \mw{कोनी।}} — \emph{dukān āj band konī.} — The shop isn't closed today.
\end{itemize}

\begin{tcolorbox}[breakable,colback=teal!4,colframe=teal!35!black,title={Before you move on}]
What does \textbf{\mw{बंद}} mean? (\textquotedblleft{}Closed\textquotedblright{}.) Does it change for a shop? (No.)
\end{tcolorbox}
\section[\texorpdfstring{\mw{तांई} (tā̃ī)}{तांई (tā̃ī)}]{\mw{तांई} (tā̃ī) — until}

\label{lesson:MW-C164-tai}



\begin{quote}
Before the new one: say \emph{he comes from the market}. Then \emph{the shop is closed}.
\end{quote}

\subsection*{You'll want to know: \mw{तांई}}
\textbf{\mw{तांई}} — \emph{tā̃ī} — \textquotedblleft{}until\textquotedblright{}.

\begin{quote}
\textbf{\mw{दुकान} \mw{दिनूगै} \mw{६} \mw{सूं} \mw{१०} \mw{तांई} \mw{खुली} \mw{है।}}
\end{quote}

\begin{quote}
\emph{dukān dinūgai chha sū̃ das tā̃ī khulī hai.} — The shop is open in the morning from six until ten.
\end{quote}

\textbf{\mw{सूं}} opens the stretch of time, and \textbf{\mw{तांई}} closes it. The hours are digits, read as numbers: \textbf{\mw{६}} is \emph{chha}, \textbf{\mw{१०}} is \emph{das}.

\begin{cousinweb}
\textbf{\mw{तांई}} (\emph{tā̃ī}) is probably inherited from Sanskrit \emph{tāvat}, \textquotedblleft{}so long\textquotedblright{}. CDIAL gives no Marwari form in this entry, but it lists an Old Gujarati word for \textquotedblleft{}up till\textquotedblright{} there (CDIAL 5804).
\end{cousinweb}

\begin{grammarlens}[title={What you've built}]
\textbf{\mw{सूं}} … \textbf{\mw{तांई}} frames any stretch: \textbf{\mw{आथण} \mw{५} \mw{सूं} \mw{८} \mw{तांई}}, in the evening from five until eight.
\end{grammarlens}

\subsection*{Your turn}
\begin{itemize}
\raggedright
  \item \emph{Say it:} \emph{tā̃ī}
  \item \emph{Say it:} \textbf{\mw{आथण} \mw{५} \mw{सूं} \mw{८} \mw{तांई}} — \emph{āthaṇ pāṅch sū̃ āṭh tā̃ī} — in the evening, from five until eight
  \item \emph{Recall:} say \emph{from the market}, then \emph{tā̃ī} again
\end{itemize}

\begin{tcolorbox}[breakable,colback=teal!4,colframe=teal!35!black,title={Before you move on}]
What does \textbf{\mw{तांई}} mean? (\textquotedblleft{}Until\textquotedblright{}.) \emph{From six until ten}? (\textbf{\mw{६} \mw{सूं} \mw{१०} \mw{तांई}}.)
\end{tcolorbox}
\section[(reading)]{(reading) — a notice at the dairy}

\label{lesson:MW-C164-dairy-notice}



\begin{quote}
Say \emph{open}, \emph{closed} and \emph{until}.

Nobody is talking to anybody here. A notice says what you may do, and when.
\end{quote}

\begin{tcolorbox}[breakable,skin=enhanced,colback=orange!6,colframe=orange!45!black,boxrule=0.5pt,arc=1mm,left=6pt,right=6pt,top=4pt,bottom=4pt,fonttitle=\bfseries,title={Read this}]
\begin{quote}
\textbf{\mw{दूध}, \mw{दही} \mw{अर} \mw{छाछ}}
\end{quote}

\begin{quote}
\textbf{\mw{दुकान} \mw{रोज} \mw{दिनूगै} \mw{६} \mw{सूं} \mw{१०} \mw{तांई} \mw{खुली} \mw{है।}} \textbf{\mw{आथण} \mw{दुकान} \mw{५} \mw{सूं} \mw{८} \mw{तांई} \mw{खुली} \mw{है।}} \textbf{\mw{रविवार} \mw{दुकान} \mw{दिनूगै} \mw{६} \mw{सूं} \mw{१२} \mw{तांई} \mw{खुली} \mw{है।}} \textbf{\mw{दोपार} \mw{अर} \mw{रात} \mw{दुकान} \mw{बंद} \mw{है।}} \textbf{\mw{सोमवार} \mw{दुकान} \mw{बंद} \mw{है।}} \textbf{\mw{भाव}: \mw{दूध} \mw{५४।}} \textbf{\mw{आज} \mw{घी} \mw{कोनी।}} \textbf{\mw{दही} \mw{घणो} \mw{आछो} \mw{है}!} \textbf{\mw{आभार।}}
\end{quote}

\emph{Read it:} the notice once, straight through, without translating

\emph{Read it:} it a second time, and find every line that says \emph{when}
\end{tcolorbox}

\subsection*{What to know first}
Most lines open on a \emph{when}: \textbf{\mw{रोज}}, \textbf{\mw{आथण}}, \textbf{\mw{रविवार}}, \textbf{\mw{दोपार}}, \textbf{\mw{सोमवार}}. Each stretch of hours runs \textbf{\mw{सूं}} … \textbf{\mw{तांई}}. \textbf{\mw{भाव}} is the price, written in digits. The line with \textbf{!} is an advert, set into the notice to bring you in.

\subsection*{Your turn}
\begin{itemize}
\raggedright
  \item \emph{Answer:} is the shop open on Monday? (No: \textbf{\mw{सोमवार} \mw{दुकान} \mw{बंद} \mw{है।}})
  \item \emph{Answer:} until when is it open on Sunday morning? (Until twelve: \textbf{\mw{१२} \mw{तांई}}.)
  \item \emph{Answer:} is there ghee today? (No: \textbf{\mw{आज} \mw{घी} \mw{कोनी।}})
\end{itemize}

\begin{tcolorbox}[breakable,colback=teal!4,colframe=teal!35!black,title={Before you move on}]
Which day is the shop closed all day? (\textbf{\mw{सोमवार}}, Monday.) When is it closed every day? (\textbf{\mw{दोपार} \mw{अर} \mw{रात}}, midday and night.)
\end{tcolorbox}
\section[(dialogue)]{A notice at the dairy — a review}

\label{lesson:MW-R164-dairy}



\begin{quote}
Say \emph{the shop is open}.
\end{quote}

\subsection*{Your turn}
\begin{itemize}
\raggedright
  \item \emph{Say it:} \emph{the house is open}; \emph{the shop is open}; \emph{the shop is closed}
  \item \emph{Say it:} \emph{from six until ten}; \emph{in the evening, from five until eight}
  \item \emph{Say it:} \emph{on Monday the shop is closed}; \emph{there is no ghee today}
  \item \emph{Answer:} in the notice, when can you buy milk on a Sunday morning? (From six until twelve.)
\end{itemize}

\begin{tcolorbox}[breakable,colback=teal!4,colframe=teal!35!black,title={Before you move on}]
\emph{Open}? (\textbf{\mw{खुलो}}, \textbf{\mw{खुली}}.) \emph{Closed}? (\textbf{\mw{बंद}}.) \emph{Until}? (\textbf{\mw{तांई}}.)
\end{tcolorbox}
